% !PRINTED 272
\weberpage{272}
%<<CONT
is, and every divisor of these numbers also occurs among the degrees of the cyclic groups. The index of a cyclic divisor can never be smaller than $1/2\,(p^2-1)$, $p(p+1)$, $p(p-1)$.

If the degree of a divisor $G$ of $L_p$ is divisible by $p$, then it contains a cyclic group of degree $p$, and by \S.\,68 we can transform the group $G$ in such a way that this cyclic group consists of the substitutions
\[
T = \begin{pmatrix}1, & t\\ 0, & 1\end{pmatrix}, \quad t = 0, 1, 2, \ldots, p-1
\]
If now $G$ contains a further substitution
\[
A = \begin{pmatrix}a, & b\\ c, & d\end{pmatrix}
\]
in which $c$ is different from zero, then
\weq{1}{\begin{pmatrix}1, & -ac^{-1}\\ 0, & 1\end{pmatrix}\begin{pmatrix}a, & b\\ c, & d\end{pmatrix}\begin{pmatrix}1, & -dc^{-1}\\ 0, & 1\end{pmatrix} = \begin{pmatrix}0, & -c^{-1}\\ c, & 0\end{pmatrix}}
also occurs in it. From this it follows further that
\[
\begin{pmatrix}0, & -c^{-1}\\ c, & 0\end{pmatrix}\begin{pmatrix}1, & t\\ 0, & 1\end{pmatrix}\begin{pmatrix}0, & -c^{-1}\\ c, & 0\end{pmatrix} = \begin{pmatrix}1, & 0\\ -c^2 t, & 1\end{pmatrix},
\]
and consequently every substitution
\[
U = \begin{pmatrix}1, & 0\\ u, & 1\end{pmatrix}, \quad u = 0, 1, 2, \ldots, p-1
\]
also occurs. If one puts $U$ for $u = -1$ in place of $A$ in the formula (1), it follows that $G$ also contains the substitution $\begin{pmatrix}0, & 1\\ -1, & 0\end{pmatrix}$, which together with $U$ gives a system of generating elements of $L_p$ (\S.\,65), so that $G$ is therefore identical with $L_p$. From this it follows that $G$, if it is not to be identical with $L_p$, can contain only substitutions of the form
\[
A = \begin{pmatrix}a, & b\\ 0, & a^{-1}\end{pmatrix}
\]
Conversely, all these substitutions form a group of degree $1/2\,p(p-1)$, and hence a divisor of index $p+1$.

Among these there are also contained certain divisors of greater index, which one obtains by taking for $a$ not all values, but all values of the form $a^s$, where $s$ is a divisor of $1/2\,(p-1)$. The degree of such a group is $1/2\,p(p-1) : s$.
